\documentclass[11pt,a4paper]{article}
\usepackage[utf8]{inputenc}
\usepackage[T1]{fontenc}
\usepackage[spanish,es-noquoting]{babel}
\usepackage[left=2.3cm,right=2.3cm,top=2.3cm,bottom=2.3cm]{geometry}
\usepackage{amsmath,amssymb}
\usepackage{booktabs}
\usepackage{tabularx}
\usepackage{enumitem}
\usepackage{titlesec}
\usepackage{xcolor}
\definecolor{gris}{gray}{0.35}
\usepackage[hidelinks]{hyperref}

\titleformat{\section}{\large\bfseries}{\thesection.}{0.5em}{}
\titlespacing*{\section}{0pt}{1.4em}{0.6em}
\titleformat{\subsection}{\normalsize\bfseries}{}{0em}{}
\titlespacing*{\subsection}{0pt}{1em}{0.35em}

\newcolumntype{Y}{>{\raggedright\arraybackslash}X}
\setlist[itemize]{topsep=3pt,itemsep=2pt,leftmargin=1.3em}
\setlist[enumerate]{topsep=3pt,itemsep=2pt,leftmargin=1.6em}
\renewcommand{\arraystretch}{1.2}

% babel-spanish redefine \% para insertar un espacio fino antes, y para decidirlo
% compara \lastskip (que dentro de tabularx viene en em) contra cero (en pt).
% Eso lanza "Incompatible glue units" y aborta. Se restaura el \% de LaTeX.
\AtBeginDocument{%
  \renewcommand{\%}{\char`\%\relax}%  \relax evita que TeX coma el espacio
  \decimalpoint                          % punto decimal, como en la tesis
}

\newcommand{\cita}[1]{%
  \begin{quote}\itshape\color{gris}#1\end{quote}}

\pagestyle{plain}

\begin{document}

\begin{center}
  {\Large\bfseries Puntos abiertos del Capítulo 1}\\[0.3em]
  {\large para revisión con el Dr. José Adán Hernández Nolasco}\\[0.9em]
  \begin{tabular}{rl}
    \textbf{Autor:} & Roberto Hernández Estrada \\
    \textbf{Fecha:} & 18 de septiembre de 2026 \\
  \end{tabular}
\end{center}

\vspace{0.6em}
\noindent\rule{\textwidth}{0.4pt}
\vspace{0.4em}

\noindent
El Capítulo 1 (Generalidades) conserva el planteamiento del protocolo de octubre
de 2025. El trabajo realizado desde entonces cambió el método, y hay
afirmaciones del capítulo que el código ya no respalda.

\noindent
\textbf{El punto 1 ---el título--- ya se resolvió} en la conversación del 18 de
septiembre y está aplicado. Los cuatro restantes siguen abiertos y se presentan
aquí con su texto actual, lo que muestran los resultados, y la decisión concreta
que hay que tomar. \textbf{Ninguno se ha modificado.}

\section*{Dónde está el trabajo hoy}

\begin{tabularx}{\textwidth}{Yll}
\toprule
\textbf{Experimento} & \textbf{Resultado} & \textbf{Estado} \\
\midrule
NB01 --- Helmholtz 1D                & $L^2 = 0.006\,\%$ & Cerrado \\
NB02 --- Helmholtz 2D, campo complejo & $L^2 = 0.171\,\%$ & Cerrado \\
NB02B --- puente de arquitectura      & $L^2 = 0.0078\,\%$ & Cerrado \\
NB03 --- speckle en $z=1\lambda$      & $L^2 = 3.161\,\%$, 5/5 & Cerrado \\
NB03B --- speckle en $z=2\lambda$     & $L^2 = 0.993\,\%$, 5/5 & Cerrado \\
NB03C --- speckle hasta $z=5\lambda$  & $L^2 = 2.244\,\%$, 5/5 & Cerrado \\
NB03D --- speckle hasta $z=10\lambda$ & $L^2 = 2.617\,\%$, 5/5 & Cerrado \\
\textbf{NB04 --- benchmark contra FEM} & --- & \textbf{No iniciado} \\
\bottomrule
\end{tabularx}

\vspace{0.7em}
\noindent
La tesis está en 81 páginas, sin errores de compilación, y 369 cifras del texto
se verifican automáticamente contra los archivos de resultados que las
generaron.

\vspace{0.4em}
\noindent
\textbf{Novedad desde la última versión de este documento:} la cadena de
subdominios se extendió a $z=10\lambda$ con diez bloques, reutilizando sin
reentrenar el tramo $0$--$5\lambda$ ya validado. Duplicar la distancia eleva el
error de $2.244\,\%$ a $2.617\,\%$ ---un incremento de $0.373$ puntos por cinco
bloques adicionales---, con las cinco pantallas aceptadas. Una cadena de veinte
bloques hasta $20\lambda$, ensayada sobre una sola pantalla y sin refinamiento,
\emph{no} alcanza el criterio ($8.743\,\%$).

\vspace{0.4em}
\noindent
\textbf{El hallazgo que condiciona todo lo demás:} la formulación directa que
planteaba el protocolo resultó estar \textbf{mal planteada}. Cumplía la ecuación
de Helmholtz a $10^{-14}$ y la frontera exactamente, y aun así daba
\textbf{92.75\,\% de error}: sin condición de radiación el problema admite
infinitas soluciones. Se sustituyó por una formulación modal con condición de
Cauchy dura, que es la regularización canónica de este problema
(Nguyen et al., 2014).

\section{El título \normalfont\textit{(resuelto el 18/09)}}

\subsection{Texto anterior}
\cita{Simulación \textbf{acelerada} de speckle óptico: \textbf{Un enfoque
basado} redes neuronales físicamente informadas (PINN's)}

\subsection{Qué dice el trabajo}
\begin{itemize}
  \item \textbf{«Acelerada» no tiene respaldo.} El factor de aceleración se
  define contra FEM, y ese benchmark no se ha hecho. Lo que sí está medido es la
  comparación contra el espectro angular: \textbf{18.5$\times$ a favor del
  espectro angular} para un plano, y \textbf{1.76$\times$ a favor de la PINN}
  para 101 planos simultáneos. Es un resultado condicional, y un título no
  admite condiciones.
  \item El \textbf{8 de septiembre} usted pidió quitar «acelerada» y «un enfoque
  basado».
  \item «(PINN's)» lleva un apóstrofo que no corresponde al plural.
\end{itemize}

\subsection{Decisión}
Propuesta, 14 palabras:

\begin{quote}
\bfseries
Simulación del speckle óptico mediante Redes Neuronales Informadas por Física:
formulación modal de Helmholtz
\end{quote}

\noindent
La parte anterior a los dos puntos es la que quedó acordada en esa llamada. El
añadido es propuesta mía: \textbf{la formulación modal es precisamente lo que
distingue el trabajo entregado del que prometía el protocolo}, y es lo que lo
separa de Panagiotakopoulos et al. (2026), que valida cualitativamente sin dar
cifras.

\vspace{0.3em}
\noindent
\textbf{Aprobado y aplicado.} El título figura ya en la portada, y se añadió al
Dr.~Noel Zacarías Morales como coasesor. Se conserva aquí el razonamiento
porque sostiene los puntos 2 y 3.

\section{El objetivo general dice «resolución directa»}

\subsection{Texto actual (Cap. 1, \textit{Objetivo general})}
\cita{\ldots para la simulación acelerada y físicamente precisa de patrones de
speckle óptico, mediante la \textbf{resolución directa} de la ecuación de
Helmholtz bidimensional con campo eléctrico complejo.}

\subsection{Qué dice el trabajo}
La resolución directa es exactamente el enfoque que \textbf{se abandonó}. El
método vigente expande el campo en serie de Fourier ---exacta, por las fronteras
laterales periódicas--- y resuelve 41 ecuaciones diferenciales ordinarias
desacopladas, una por modo propagante.

El objetivo general, tal como está, describe un método que la tesis ya no usa.

\subsection{Decisión}
Propuesta:

\cita{\ldots mediante una \textbf{formulación modal} de la ecuación de Helmholtz
bidimensional con campo eléctrico complejo, con la condición de Cauchy impuesta
por construcción.}

\noindent
\textbf{¿Se corrige, o se prefiere otra redacción?}

\section{La hipótesis (ii) no tiene respuesta}

\subsection{Texto actual}
\cita{(ii) obtener un factor de aceleración $S = T_{\mathrm{FEM}} /
T_{\mathrm{PINN}} > 1$ en la fase de inferencia}

\subsection{Qué dice el trabajo}
NB04 no se ha hecho, así que \textbf{la hipótesis (ii) queda sin contestar}.
Tres obstáculos concretos:

\begin{enumerate}
  \item \textbf{FEniCSx no corre en Windows nativo.} Requiere WSL2 o Docker, no
  instalados.
  \item \textbf{La comparación está condicionada} por el propio Capítulo 1 a que
  el problema se formule con fronteras absorbentes (ver punto 4).
  \item \textbf{El resultado sería probablemente desfavorable.} El entrenamiento
  cuesta 351.9 s por pantalla, y el espectro angular ya resultó más rápido en
  evaluación de un solo plano.
\end{enumerate}

\noindent
Conviene señalar que el protocolo de octubre pedía \textbf{una aceleración de al
menos 10$\times$}; el texto actual de la tesis la relajó a «mayor que 1». Esa
rebaja está en el documento y conviene decidirla explícitamente.

\subsection{Decisión --- tres opciones}

\begin{tabularx}{\textwidth}{lY}
\toprule
\textbf{Opción} & \textbf{Qué implica} \\
\midrule
\textbf{A.} Hacer NB04 & Instalar WSL2 o Docker. Cierra la hipótesis, sea cual
sea el resultado. Requiere resolver antes el punto 4 \\
\textbf{B.} Reformular & Cambiarla por el benchmark que sí está hecho:
evaluación libre de malla frente a una referencia de propagación. Honesto, y ya
está medido \\
\textbf{C.} Trabajo futuro & Dejarla enunciada y explicar en Cap. 5 por qué no
se resolvió \\
\bottomrule
\end{tabularx}

\vspace{0.5em}
\noindent
Mi lectura: la contribución de la tesis quedó en \textbf{precisión y buen
planteamiento}, no en velocidad. El Capítulo 4 ya está escrito en esos términos.

\subsection{Existe un precedente, y es de esta misma dirección}

La tesis de Gabriel Omar Domínguez Ruiz, aprobada el 9 de julio de 2026 bajo su
dirección, se topó con esta misma cuestión y la resolvió. Su Sección~3.4
compara contra métodos clásicos y \textbf{el resultado le es desfavorable}:
Runge--Kutta alcanza errores de $10^{-9}$ en milisegundos y las diferencias
finitas llegan a $10^{-6}$, frente al $0.2\,\%$ de sus PINN. Su redacción:

\cita{…la aportación de las PINN \textbf{no debe entenderse como una
superioridad en exactitud puntual o en velocidad de cálculo}, sino en la
\emph{naturaleza} de la solución que cada enfoque produce.}

\cita{La comparación \textbf{no es una competencia por la menor cifra de
error} ---en una sola corrida, los métodos numéricos son rápidos y muy
precisos---, sino el reconocimiento de \textbf{dos objetos distintos}.}

Desplazó la afirmación de velocidad a diferenciabilidad y ausencia de malla, y
lo acompañó de una tabla comparativa cualitativa. Ese planteamiento ya pasó por
el comité.

\vspace{0.4em}
\noindent
Este trabajo está además en mejor posición que aquél en un punto: \textbf{el
benchmark contra la referencia de propagación ya está hecho}, y arroja
$1.76\times$ a favor de la PINN en evaluación de $101$ planos simultáneos. No se
partiría de cero.

\vspace{0.3em}
\noindent
\textbf{¿Qué opción se toma?}

\section{Fronteras absorbentes frente a periódicas}

\subsection{Texto actual (Cap. 1, \textit{Alcances})}
\cita{Se realizará una evaluación comparativa del tiempo de inferencia contra
una referencia numérica de propagación y, \textbf{cuando el problema se formule
con fronteras absorbentes}, contra una implementación FEM.}

\subsection{Qué dice el trabajo}
NB03 usa \textbf{fronteras laterales periódicas}, y no por comodidad: la
periodicidad es lo que hace que la serie de Fourier sea \textbf{exacta} y que
los 41 modos queden desacoplados. Cambiarlas a absorbentes destruiría la
formulación modal entera.

Es decir: \textbf{la condición que el Capítulo 1 pone para comparar contra FEM
es incompatible con el método que la tesis usa.}

\subsection{Decisión}
\begin{itemize}
  \item \textbf{A.} Retirar la condición, y comparar contra FEM con fronteras
  periódicas
  \item \textbf{B.} Mantenerla, aceptando que NB04 exige reformular el problema
  \item \textbf{C.} Retirar la mención a FEM del alcance, coherente con el
  punto 3
\end{itemize}

\noindent
\textbf{¿Cuál?}

\section*{Una oportunidad, no un problema}

El protocolo incluía una tercera pregunta de investigación que la tesis ya no
plantea:

\cita{¿Cómo impacta la complejidad del medio difusor (rugosidad de la
superficie) en la precisión y el tiempo de entrenamiento del modelo PINN?}

\noindent
\textbf{Hay material para contestarla.} Existen en disco referencias generadas
con cinco longitudes de correlación distintas ---0.10, 0.20, 0.30, 0.40 y
0.50\,$\lambda$---, y la tesis sólo reporta la de 0.10\,$\lambda$.

\noindent
Recuperar esa pregunta costaría un barrido comparativo, sin formulación nueva ni
arquitectura nueva. \textbf{¿Vale la pena, o se deja fuera del alcance?}

\section*{Resumen de decisiones}

\begin{tabularx}{\textwidth}{clY}
\toprule
\textbf{\#} & \textbf{Punto} & \textbf{Decisión pedida} \\
\midrule
1 & Título & \textit{Resuelto el 18/09: aprobado y aplicado} \\
2 & Objetivo general & ¿Se corrige «resolución directa» $\rightarrow$
«formulación modal»? \\
3 & Hipótesis (ii) & ¿Hacer NB04, reformular, o declarar trabajo futuro? \\
4 & Fronteras absorbentes & ¿Retirar la condición, mantenerla, o quitar FEM del
alcance? \\
5 & Rugosidad & ¿Se recupera la pregunta del protocolo? \\
\bottomrule
\end{tabularx}

\vspace{0.6em}
\noindent
Los puntos \textbf{2 y 3 están ligados}: si la hipótesis (ii) se reformula, el
objetivo general debe decir lo mismo. Y el \textbf{4 condiciona al 3}: no se
puede decidir sobre NB04 sin resolver antes la cuestión de las fronteras.

\end{document}
